%! Adding \& Subtracting Rational Expressions — Unit 4, Lesson 4.3 — Guided Notes (Answer Key)
\documentclass[10pt]{article}
\usepackage{algebra2-article}
\usepackage{algebra2-key}   % pulls in -boxes; do NOT also load -boxes
\newcommand{\TallMath}[1]{$\displaystyle #1\rule[-1.4em]{0pt}{3.2em}$}
\newcommand{\lgrid}[4]{%
  \draw[step=1, linegray!40, very thin] (#1,#3) grid (#2,#4);
  \draw[->, semithick, charcoal] (#1,0)--(#2,0) node[below right, font=\scriptsize]{$x$};
  \draw[->, semithick, charcoal] (0,#3)--(0,#4) node[left, font=\scriptsize]{$y$};}
% Vocab answer for the key: bold forest term, then the filled definition on a write-line.
% The leading and trailing \par are required: \ansline ends with \dotfill but does NOT end the
% paragraph, so without them each term label gets pulled onto the previous answer's line.
\newcommand{\vocabans}[2]{%
  \par\noindent\textbf{\textcolor{forest}{#1:}}\\[1pt]\ansline{#2}\par}

\begin{document}

\pageheader{Unit 4, Lesson 4.3}{Guided Notes --- Answer Key}
\namedateperiod

\vspace{0.08in}

\begin{objectivebox}
\textbf{By the end of this lesson, I will be able to\dots}
\begin{itemize}\setlength{\itemsep}{2pt}
  \item add and subtract rational expressions with like denominators, putting the second numerator in \emph{parentheses} whenever I subtract;
  \item build the \textbf{least common denominator} from the \emph{factored} denominators, rewrite each fraction to match it, combine, and simplify;
  \item recognize a \textbf{complex fraction} and simplify it as a quotient of two simple fractions --- and state every restriction it carries.
\end{itemize}
\end{objectivebox}

\vspace{0.05in}

\begin{vocabbox}
\small
Fill in each term as we build it together.
\par\vspace{2pt}   % \par is required: \vocabans's \noindent is a no-op mid-paragraph
\vocabans{Least common denominator (LCD)}{the smallest expression every denominator divides into --- built by taking each \emph{distinct} factor to the \emph{highest} power it reaches in any one denominator}
\vocabans{Building factor}{the form of $1$, written $\frac{k}{k}$, that you multiply a fraction by to give its denominator whatever the LCD has and it is missing}
\vocabans{Complex algebraic fraction}{a fraction that has fractions inside its numerator, its denominator, or both --- simplified as a quotient of two simple fractions}
\vocabans{Main fraction bar}{the long bar of a complex fraction; it is a division sign, so whatever sits below it may never equal zero}
\end{vocabbox}

\begin{hookbox}
Two students subtracted the same expression and got different answers. Let
\[
  D(x)=\frac{4x-1}{x+3}-\frac{2x-7}{x+3} .
\]
\begin{center}
\renewcommand{\arraystretch}{1.35}
\begin{tabular}{|l|c|c|c|}
\hline
 & $x=0$ & $x=1$ & $x=-3$ \\ \hline
$D(x)$ \; (the original) & \ans{$2$} & \ans{$2$} & \ans{undef.} \\ \hline
Dev: \; $\dfrac{2x-8}{x+3}$ & \ans{$-\frac83$} & \ans{$-\frac32$} & \ans{undef.} \\ \hline
Elin: \; $2$ & \ans{$2$} & \ans{$2$} & \ans{$2$} \\ \hline
\end{tabular}
\end{center}
\begin{itemize}\setlength{\itemsep}{1pt}
  \item Dev wrote $\dfrac{4x-1-2x-7}{x+3}$. Which \emph{one} keystroke is missing, and what did losing
        it do to the sign of the $7$? \ansline{the parentheses around $2x-7$. Subtracting the whole numerator makes it $-(2x-7)=-2x+7$; Dev kept the $-7$, so his numerator is off by $14$.}
  \item Elin is right about the value --- but look at the last column. What must she add to her answer
        before it is finished? \ansline{the restriction $x\neq-3$. The original has no value at $x=-3$, so the answer is ``$2$, for every $x$ except $-3$'' --- a horizontal line with a hole in it.}
\end{itemize}
Two ideas run the whole period. \textbf{A minus sign in front of a fraction belongs to that fraction's
entire numerator.} And once again: even an answer as plain as ``$2$'' still carries the restrictions of
the problem it came from.
\end{hookbox}

\begin{notesbox}{1. Like Denominators: Add the Numerators, Keep the Denominator}
When the denominators already match, there is nothing to build:
\[
  \frac{a}{c}+\frac{b}{c}=\frac{a+b}{c}, \qquad
  \frac{a}{c}-\frac{b}{c}=\frac{a-b}{c}, \qquad c\neq0 .
\]

\vspace{3pt}
\textbf{You never add the denominators.} Check it on numbers you already know:
$\dfrac12+\dfrac12=$ \ans{1}, \; but adding denominators would give
$\dfrac{2}{4}=$ \ans{$\frac12$}. Adding denominators would make two halves into
\ans{one half}, which is nonsense.

\vspace{4pt}
\textbf{Worked.} \; $\dfrac{7}{x^2-4}+\dfrac{x}{x^2-4}$
\begin{itemize}[leftmargin=1.5em, itemsep=3pt]
  \item Factor the denominator \emph{first}, so you can see the restrictions: $x^2-4=$ \ans{$(x-2)(x+2)$},
        so $x\neq$ \ans{$2,\ -2$}.
  \item Add the numerators, keep the denominator: \ans{$\frac{x+7}{(x-2)(x+2)}$}
  \item \textbf{Last step, always:} factor the new numerator and check whether anything divides out.
        Here $x+7$ is prime, so we are done.
\end{itemize}
\end{notesbox}

\begin{notesbox}{2. Subtracting: the Minus Sign Owns the Whole Numerator}
\begin{center}
\fbox{\parbox{0.90\linewidth}{\centering\vspace{2pt}
\textbf{Put the second numerator in parentheses before you subtract.} Then distribute the minus sign
--- it changes the sign of \emph{every} term inside.\vspace{2pt}}}
\end{center}

\vspace{4pt}
\textbf{Worked.} \; $\dfrac{6x+5}{x^2-4}-\dfrac{2x+13}{x^2-4}$
\begin{itemize}[leftmargin=1.5em, itemsep=3pt]
  \item Write the subtraction with parentheses:
        $\dfrac{(6x+5)-({\color{keyred}\mathbf{2x+13}})}{x^2-4}$
  \item Distribute and collect: numerator $=$ \ans{$4x-8$} $=$ \ans{$4(x-2)$} (factored).
  \item Divide out, then attach the restrictions from the \emph{original} denominator:
        \ans{$\frac{4}{x+2}$}, \; $x\neq$ \ans{$2,\ -2$}.
  \item Check it at $x=0$: the original gives \ans{2} and your answer gives \ans{2}.
\end{itemize}

\vspace{3pt}
Notice which restriction survived. The factor $(x-2)$ divided out, so $x=2$ is a
\ans{hole} in the graph; $(x+2)$ stayed, so $x=-2$ is a \ans{wall}. That sort is Lesson 4.0's,
and it never stops applying.
\end{notesbox}

\begin{notesbox}{3. Unlike Denominators: Build the LCD}
You cannot combine pieces that are cut differently. The \textbf{least common denominator} is the
smallest expression every denominator divides into --- and you build it out of \emph{factors}, never by
multiplying the denominators together.

\vspace{4pt}
\textbf{The six steps.}
\begin{enumerate}[leftmargin=1.9em, itemsep=2pt]
  \item \textbf{Factor} every denominator completely.
  \item \textbf{Build the LCD:} take every \emph{distinct} factor, each to the \emph{highest} power it reaches in any one denominator.
  \item \textbf{List the restrictions} --- set each factor of the LCD equal to zero.
  \item \textbf{Rewrite} each fraction: multiply it by the \textbf{building factor} $\frac{k}{k}$ that supplies whatever its denominator is missing.
  \item \textbf{Combine} the numerators over the LCD (parentheses on every subtraction) and simplify.
  \item \textbf{Factor the numerator} and divide out if you can. Attach the restrictions.
\end{enumerate}

\vspace{4pt}
\textbf{Monomial denominators.} \; $\dfrac{5}{6x^2}+\dfrac{7}{4x}$
\begin{itemize}[leftmargin=1.5em, itemsep=3pt]
  \item $6x^2=2\cdot3\cdot x^2$ and $4x=2^2\cdot x$, so LCD $=$ \ans{$12x^2$}.
  \item Building factors: the first needs \ans{$\frac22$}, the second needs \ans{$\frac{3x}{3x}$}.
  \item Result: \ans{$\frac{21x+10}{12x^2}$}, \; $x\neq$ \ans{0}.
\end{itemize}

\vspace{5pt}
\textbf{The anchor.} \; $\dfrac{3}{x-3}-\dfrac{18}{x^2-9}$
\begin{itemize}[leftmargin=1.5em, itemsep=3pt]
  \item Factor: $x^2-9=$ \ans{$(x-3)(x+3)$}, so LCD $=$ \ans{$(x-3)(x+3)$} and $x\neq$ \ans{$3,\ -3$}.
  \item The first fraction is missing the factor \ans{$x+3$}; multiply it by that over itself:
        $\dfrac{3}{x-3}\cdot\dfrac{{\color{keyred}\mathbf{x+3}}}{{\color{keyred}\mathbf{x+3}}}=
        \dfrac{{\color{keyred}\mathbf{3x+9}}}{(x-3)(x+3)}$
  \item Combine and simplify the numerator: \ans{$3x-9$} $=$ \ans{$3(x-3)$} (factored).
  \item Divide out. Final answer: \ans{$\frac{3}{x+3}$}, \; $x\neq$ \ans{$3,\ -3$}.
\end{itemize}
Where have you seen that answer before today? \ansline{it is exactly the Warm-Up's item 3 --- the simplified form of $\frac{3x-9}{x^2-9}$, restrictions and all. Two different starting points, one function.}

\vspace{5pt}
\textbf{Two trinomial denominators.} \; $\dfrac{2}{x^2+5x+6}+\dfrac{3}{x^2-9}$
\begin{itemize}[leftmargin=1.5em, itemsep=3pt]
  \item Factored: \ans{$(x+2)(x+3)$} and \ans{$(x-3)(x+3)$}. The shared factor is \ans{$x+3$}, so the LCD uses
        it \emph{once}: LCD $=$ \ans{$(x+2)(x+3)(x-3)$}.
  \item Restrictions: $x\neq$ \ans{$-2,\ -3,\ 3$}
  \item Combine: $\dfrac{2({\color{keyred}\mathbf{x-3}})+3({\color{keyred}\mathbf{x+2}})}{\text{LCD}}=$
        \ans{$\frac{5x}{(x+2)(x+3)(x-3)}$}
\end{itemize}

\vspace{4pt}
\begin{center}
\fbox{\parbox{0.92\linewidth}{\centering\vspace{2pt}
\textbf{The LCD \emph{is} the restriction list.} Every factor of every original denominator is in
there, so setting the LCD's factors to zero gives you every excluded value --- and it costs you no
extra work.\vspace{2pt}}}
\end{center}
\vspace{2pt}
One warning attached to that shortcut: build the LCD from the \emph{original} factored denominators.
Reading anything off your final answer is the error from Lesson 4.1, and it is still wrong today.

\vspace{5pt}
\textbf{Opposite binomials in the denominators} (Lesson 4.1, still in force). \;
$\dfrac{4}{x-5}+\dfrac{2}{5-x}$
\begin{itemize}[leftmargin=1.5em, itemsep=3pt]
  \item These denominators are not different --- they are opposites: $5-x=$ \ans{$-1\cdot$} $(x-5)$.
  \item Rewrite the second fraction as $\dfrac{{\color{keyred}\mathbf{-2}}}{x-5}$. Now the denominators match.
  \item Answer: \ans{$\frac{2}{x-5}$}, \; $x\neq$ \ans{5}. \; (The LCD was never $(x-5)(5-x)$.)
\end{itemize}
\end{notesbox}

\begin{notesbox}{4. Complex Fractions: a Division in Disguise}
A \textbf{complex fraction} has fractions inside its numerator, its denominator, or both. The long
\textbf{main fraction bar} is not decoration --- it is a division sign.
\[
  \frac{\;\dfrac{1}{x}+\dfrac{1}{2}\;}{\;\dfrac{1}{x}-\dfrac{1}{4}\;}
  \qquad\text{means}\qquad
  \left(\frac{1}{x}+\frac{1}{2}\right)\div\left(\frac{1}{x}-\frac{1}{4}\right)
\]

\vspace{3pt}
\textbf{Method 1 --- combine, then divide.} Turn the top into one simple fraction, turn the bottom into
one simple fraction, then keep--change--flip (Lesson 4.2).
\begin{itemize}[leftmargin=1.5em, itemsep=3pt]
  \item Top: $\dfrac{1}{x}+\dfrac{1}{2}=$ \ans{$\frac{x+2}{2x}$} \qquad Bottom: $\dfrac{1}{x}-\dfrac{1}{4}=$
        \ans{$\frac{4-x}{4x}$}
  \item Divide: \; \ans{$\frac{x+2}{2x}$} $\cdot$ \ans{$\frac{4x}{4-x}$} $=$ \ans{$\frac{2(x+2)}{4-x}$}
\end{itemize}

\vspace{3pt}
\textbf{Method 2 --- clear all the little fractions at once.} Multiply the top \emph{and} the bottom by
the LCD of every small fraction inside. Here that LCD is \ans{$4x$}.
\begin{itemize}[leftmargin=1.5em, itemsep=3pt]
  \item Top becomes \ans{$4+2x$}; \; bottom becomes \ans{$4-x$}; \; so the answer is
        \ans{$\frac{2(x+2)}{4-x}$} --- the same one.
\end{itemize}

\vspace{4pt}
\textbf{Restrictions: two levels, and the second one hides.}
\begin{center}
\renewcommand{\arraystretch}{1.5}
\begin{tabularx}{\linewidth}{@{}c l X@{}}
\hline
\textbf{\#} & \textbf{What must be nonzero} & \textbf{Why --- and where it hides} \\
\hline
1 & every \emph{inner} denominator & each little fraction has to exist. Visible: just read them. Here $x\neq$ \ans{0}. \\
2 & the \emph{whole bottom} & the main bar is a $\div$, and you cannot divide by zero. This is Lesson 4.2's \textbf{source 3} in a new costume: set the bottom's combined \emph{numerator} to zero. Here $4-x=0$ gives $x\neq$ \ans{4}. \\
\hline
\end{tabularx}
\end{center}
Complete restriction list for the anchor: $x\neq$ \ans{$0,\ 4$}. Check the whole thing at $x=2$: the
original gives \ans{4} and your answer gives \ans{4}.

\vspace{5pt}
\textbf{One more, with a binomial underneath.} \;
$\dfrac{\;\dfrac{1}{x}-\dfrac{1}{3}\;}{\;x-3\;}$
\begin{itemize}[leftmargin=1.5em, itemsep=3pt]
  \item Top: \ans{$\frac{3-x}{3x}$}. \; The bottom is already one piece, so divide by it:
        \ans{$\frac{3-x}{3x}$} $\cdot\;\dfrac{1}{x-3}$
  \item Watch the \emph{opposite binomials}: $3-x=$ \ans{$-1\cdot$} $(x-3)$, so the answer is
        \ans{$-\frac{1}{3x}$}.
  \item Restrictions: $x\neq$ \ans{$0,\ 3$} --- one from the inner denominator, one from the main bar.
\end{itemize}
\end{notesbox}

\begin{notesbox}{5. The Picture: One Subtraction, One Hole and One Wall}
The anchor from Section 3 was $\;g(x)=\dfrac{3}{x-3}-\dfrac{18}{x^2-9}=\dfrac{3}{x+3}$, with
$x\neq3,-3$.
\begin{center}
\begin{tikzpicture}[scale=0.42]
  \lgrid{-8}{5}{-5}{5}
  \draw[dashed, redacc, thick] (-3,-5)--(-3,5);
  \node[redacc, font=\scriptsize, above left] at (-3.1,3.4) {$x=-3$};
  \begin{scope}
    \clip (-8,-5) rectangle (5,5);
    \draw[very thick, forest, domain=-8:-3.62, samples=100, smooth]
      plot (\x,{3/((\x)+3)});
    \draw[very thick, forest, domain=-2.38:5, samples=100, smooth]
      plot (\x,{3/((\x)+3)});
  \end{scope}
  \draw[forest, very thick, fill=white] (3,0.5) circle (4pt);
  \node[charcoal, font=\scriptsize, above right] at (3.1,0.7) {hole};
\end{tikzpicture}
\end{center}
\begin{enumerate}[leftmargin=1.7em, itemsep=3pt]
  \item Both excluded values came from the \emph{original} denominators. Write those two denominators
        in factored form: \ans{$x-3$} and \ans{$(x-3)(x+3)$}. \; Which excluded value appears in
        \emph{both} of them? \ans{$x=3$}
  \item $x=3$ is a \ans{hole} because its factor \ans{divided out}; \; $x=-3$ is a
        \ans{wall} because its factor \ans{stayed downstairs}.
  \item The hole's height comes from the \emph{simplified} form: $\dfrac{3}{3+3}=$ \ans{$\frac12$}, so
        the hole sits at \ans{$\left(3,\frac12\right)$}.
  \item Your final answer $\dfrac{3}{x+3}$ shows you only one of the two restrictions. Explain, in one
        sentence, why the other one still belongs to it. \ansline{because $g$ is the \emph{original} subtraction, and at $x=3$ both of its fractions are undefined --- the simplified form only tells you what height the graph was heading for.}
\end{enumerate}
\end{notesbox}

\begin{practicebox}
Simplify completely and state \emph{all} restrictions. Show the LCD you used.

\vspace{3pt}
\begin{enumerate}[leftmargin=1.7em, itemsep=8pt]
  \item $\dfrac{x}{x+4}+\dfrac{4}{x+4}$ \hfill Simplest form: \ans{$1$} \quad
        $x\neq$ \ans{$-4$} \\[3pt]
        Your answer is a plain number. What is left of the fraction it came from? \ans{only the restriction --- a hole at $(-4,1)$}
  \item $\dfrac{5}{x}-\dfrac{2}{x+1}$ \hfill LCD: \ans{$x(x+1)$} \\[4pt]
        Simplest form: \ans{$\frac{3x+5}{x(x+1)}$} \quad $x\neq$ \ans{$0,\ -1$}
  \item $\dfrac{\;\dfrac{2}{x}+3\;}{\;\dfrac{2}{x}-3\;}$ \hfill Simplest form: \ans{$\frac{2+3x}{2-3x}$} \quad
        $x\neq$ \ans{$0,\ \frac23$} \\[3pt]
        Which of your two restrictions came from the \emph{main} fraction bar? \ans{$x=\frac23$}
\end{enumerate}
\end{practicebox}

\begin{teachernote}
\textbf{Pacing.} Sections 1--2 are fast if the Hook lands; \textbf{Section 3 is the lesson} and
\textbf{Section 4 is the standard's other half} (A2.EO.1c). If time runs short, move the
opposite-binomial example at the end of Section 3 to homework --- it is already there --- rather than
cutting Section 4 or 5.

\vspace{3pt}
\textbf{The three sentences to repeat.} \emph{The minus sign owns the whole numerator.} \emph{Build the
LCD out of factors, not by multiplying denominators.} \emph{The main bar is a division sign.} Every
error today is one of those three going unsaid.

\vspace{3pt}
\textbf{Errors to expect.} (1) Dropping the parentheses on a subtraction --- the signature error, and
worth more class time than it looks, because the wrong answer often \emph{simplifies more prettily}
than the right one (see the homework extension). Send students to Warm-Up 2(a), not to a rule.
(2) Adding the denominators, $\frac1x+\frac13=\frac{2}{x+3}$ --- kill it numerically at $x=1$, never by
assertion; this is the Tier A error analysis. (3) Using the \emph{product} of the denominators, which
is not wrong, only expensive; let a student do it and then re-factor the result in front of the class.
(4) Cancelling a term out of a sum before it is combined (Lesson 4.1's ghost). (5) Forgetting the
restriction on a constant or polynomial answer --- the Hook and Practice 1 both plant it.

\vspace{3pt}
\textbf{Complex fractions.} The standard (A2.EO.1c) literally asks students to simplify one \emph{as a
product or quotient of simple algebraic fractions}, so Method 1 is the assessed method and Method 2 is
the speed trick. Teach both, but make the point that Method 2 \emph{destroys the evidence} for
restriction source 1: multiplying through by $4x$ clears every inner denominator, so nothing on the
page remembers that $x=0$ was ever a problem. Method 1 keeps both sources readable --- the inner
denominators are still there, and the bottom's numerator is the divisor's numerator from 4.2. Ask which
version they would rather hand a grader.

\vspace{3pt}
\textbf{Section 5} is worth protecting: one subtraction produced one hole and one wall, and both came
from denominators that no longer appear. It is the same picture 4.5 will build from an equation, and
the same fact 4.6 will call an extraneous solution.
\end{teachernote}

\end{document}
